\documentclass[a4paper,12pt]{article}
\usepackage[utf8]{inputenc}
\usepackage[T1]{fontenc}
\usepackage[ngerman]{babel}
\usepackage{amsmath}
\usepackage{amssymb}
\usepackage{enumitem}
\usepackage{geometry}
\geometry{a4paper, margin=2.5cm}

\title{Einsendeaufgaben -- Aufgabe 1.2\\[0.5em]
\large 61211 Analysis}
\author{}
\date{}

\begin{document}
\maketitle

\section*{Aufgabe 1.2}

\begin{enumerate}[label=\alph*)]

    \item
    \begin{enumerate}[label=(\roman*)]

        \item
        \underline{1. $\forall x \in \mathbb{R}^n : \|x\| \geq 0$}

        \begin{align*}
            \|x\| &= |x_1| + \sum_{j=2}^{n} |x_j - x_{j-1}| \\
                  &\geq \left|x_1 + \sum_{j=2}^{n} x_j - x_{j-1}\right| \\
                  &\geq 0
        \end{align*}

        \underline{2. $\forall x \in \mathbb{R}^n : \|x\| = 0 \Longleftrightarrow x = 0$}

        \glqq$\Rightarrow$\grqq-Richtung:

        Sei $\|x\| = 0$.

        Beweis durch vollständige Induktion:

        \underline{Induktionsanfang: $n = 2$}

        Da $|x_1| + |x_2 - x_1| = 0$ gilt, muss $x_1 = 0$
        und damit $|x_2 - x_1| = |x_2| = 0$ sein.
        Also $x = \begin{pmatrix}x_1\\x_2\end{pmatrix} = 0$.

        \underline{Induktionsschritt: $n \to n+1$}

        Nach der Induktionsvoraussetzung $x_k = 0$ für alle $1 \leq k \leq n$ gilt:
        \begin{align*}
            \|x\| &= |x_1| + \sum_{j=2}^{n+1} |x_j - x_{j-1}| \\
                  &= |x_{n+1}|
        \end{align*}
        Da $\|x\| = 0$ gilt, muss $x_{n+1} = 0$ sein.
        Also $x = \begin{pmatrix}x_1\\\vdots\\x_{n+1}\end{pmatrix} = 0$.

        \glqq$\Leftarrow$\grqq-Richtung:

        Sei $x = 0$. Daraus folgt direkt
        \[
            \|x\| = |x_1| + \sum_{j=2}^{n} |x_j - x_{j-1}| = 0.
        \]

        \item
        $\forall x \in \mathbb{R}^n\; \forall a \in \mathbb{R} : \|ax\| = |a|\,\|x\|$
        \begin{align*}
            \|ax\| &= |ax_1| + \sum_{j=2}^{n} |ax_j - ax_{j-1}| \\
                   &= |a|\,|x_1| + \sum_{j=2}^{n} |a|\,|x_j - x_{j-1}| \\
                   &= |a|\left(|x_1| + \sum_{j=2}^{n} |x_j - x_{j-1}|\right) \\
                   &= |a|\,\|x\|
        \end{align*}

        \item
        $\forall x, y \in \mathbb{R}^n : \|x + y\| \leq \|x\| + \|y\|$
        \begin{align*}
            \|x + y\| &= |x_1 + y_1| + \sum_{j=2}^{n} \bigl|(x_j + y_j) - (x_{j-1} + y_{j-1})\bigr| \\
                      &= |x_1 + y_1| + \sum_{j=2}^{n} \bigl|(x_j - x_{j-1}) + (y_j - y_{j-1})\bigr| \\
                      &\leq |x_1| + \sum_{j=2}^{n} |x_j - x_{j-1}| + |y_1| + \sum_{j=2}^{n} |y_j - y_{j-1}| \\
                      &= \|x\| + \|y\|
        \end{align*}

    \end{enumerate}

    \item
    Wir setzen $\gamma = \frac{1}{n}$.
    Wir definieren $x_m = \max\{|x_1|, \ldots, |x_n|\}$.

    \underline{Fall 1: $m = 1$, $x_m = x_1$}
    \begin{align*}
        \|x\| &= |x_1| + \sum_{j=2}^{n} |x_j - x_{j-1}| \\
              &\geq |x_1| \\
              &= \frac{1}{n} |n \cdot x_1| \\
              &\geq \frac{1}{n} \sqrt{n \, x_1^2} \\
              &\geq \frac{1}{n} \sqrt{\sum_{j=1}^{n} x_j^2} \\
              &= \gamma \|x\|_2
    \end{align*}

    \underline{Fall 2: $2 \leq m \leq n$, $x_m = \max\{|x_1|, \ldots, |x_n|\}$}

    Es gilt
    \begin{align*}
        |x_1| + \sum_{j=2}^{m} |x_j - x_{j-1}|
        &= |x_m - x_{m-1}| + \cdots + |x_2 - x_1| + |x_1| \\
        &\geq |x_m - x_{m-1} + x_{m-1} - x_{m-2} + \cdots - x_2 + x_2 - x_1| + |x_1| \\
        &= |x_m - x_1| + |x_1| \\
        &\geq |x_m|
    \end{align*}

    Mit $|x_1| + \sum_{j=2}^{m} |x_j - x_{j-1}| \geq |x_m|$ können wir nun zeigen:
    \begin{align*}
        \|x\| &= |x_1| + \sum_{j=2}^{n} |x_j - x_{j-1}| \\
              &\geq |x_1| + \sum_{j=2}^{m} |x_j - x_{j-1}| \\
              &\geq |x_m| \\
              &\geq \frac{1}{n} \sqrt{n \, x_m^2} \\
              &\geq \gamma \|x\|_2
    \end{align*}

\end{enumerate}

\end{document}
